\documentclass[12pt,a4paper]{article}

\usepackage[utf8]{inputenc}
\usepackage[T1]{fontenc}
\usepackage[brazil]{babel}
\usepackage{geometry}
\usepackage{graphicx}
\usepackage{xcolor}
\usepackage{setspace}
\usepackage{enumitem}
\usepackage{newtxtext,newtxmath}
\usepackage{etoolbox}
\usepackage{multicol}
\usepackage{tabularx}

% Toggle para o Gabarito
\newtoggle{gabarito}
\settoggle{gabarito}{true} % Mude para true para ver as respostas

\geometry{
    top=1.5cm,
    bottom=2cm,
    left=1.5cm,
    right=1.5cm
}

\definecolor{corGabarito}{RGB}{180,50,50}

\newcommand{\correto}[1]{\iftoggle{gabarito}{\textcolor{corGabarito}{\textbf{#1}}}{#1}}
\newcommand{\justificativa}[1]{\iftoggle{gabarito}{\vspace{0.1cm}\noindent\textcolor{corGabarito}{\textbf{Justificativa:} #1}}{\vspace{0.6cm}}}
\newcommand{\resp}[1]{\iftoggle{gabarito}{\textcolor{corGabarito}{\textbf{#1}}}{\phantom{#1}}}

\begin{document}

\noindent
\makebox[\textwidth]{Aluno(a): \dotfill Data: \_\_/\_\_/\_\_\_\_}

\vspace{0.3cm}
\begin{center}
    \textbf{\Large Prova Teórica de Tópicos Avançados em Informática}\\
    \textbf{Técnico de Informática 2026-1}
\end{center}

\noindent
\textbf{Orientações:} Prova individual. Responda com clareza. Nas questões de Verdadeiro ou Falso, justifique as falsas.

\vspace{0.3cm}
\noindent
\textbf{Gabarito.} Preencha cada célula com a letra da opção correta, ou V/F/Sequência.

\vspace{0.2cm}
\begin{center}
\renewcommand{\arraystretch}{1.5}
\begin{tabular}{|*{15}{c|}}
\hline
01 & 02 & 03 & 04 & 05 & 06 & 07 & 08 & 09 & 10 & 11 & 12 & \hspace{0.4cm}13\hspace{0.4cm} & \hspace{0.4cm}14\hspace{0.4cm} & 15 \\
\hline
\iftoggle{gabarito}{
 \textcolor{corGabarito}{C} & \textcolor{corGabarito}{A} & \textcolor{corGabarito}{B} & \textcolor{corGabarito}{B} & \textcolor{corGabarito}{B} & \textcolor{corGabarito}{C} & \textcolor{corGabarito}{V} & \textcolor{corGabarito}{B} & \textcolor{corGabarito}{F} & \textcolor{corGabarito}{F} & \textcolor{corGabarito}{V} & \textcolor{corGabarito}{D} & \textcolor{corGabarito}{2134} & \textcolor{corGabarito}{2314} & \textcolor{corGabarito}{C} 
}{
 & & & & & & & & & & & & & & 
} \\
\hline
\end{tabular}
\end{center}
\vspace{0.1cm}
\noindent\small\textit{* Questões 13 e 14 são de relacionar colunas. Questão 15 é descritiva.}

\vspace{0.4cm}

\begin{multicols}{2}

\textbf{1.} Qual comando é utilizado para enviar definitivamente as suas modificações de código (que já foram commitadas localmente) para o repositório na nuvem (GitHub)?
\begin{enumerate}[label=(\alph*)]
    \item \texttt{git clone}
    \item \texttt{git add .}
    \item \correto{\texttt{git push origin main}}
    \item \texttt{git commit -m "atualização"}
\end{enumerate}

\vspace{0.2cm}
\textbf{2.} No paradigma de Orientação a Objetos utilizando a linguagem Python, qual a finalidade do método especial chamado \texttt{\_\_init\_\_}?
\begin{enumerate}[label=(\alph*)]
    \item \correto{Funcionar como o construtor da classe, sendo executado automaticamente para inicializar os atributos.}
    \item Importar bibliotecas externas instaladas pelo \texttt{pip}.
    \item Encerrar a conexão com o banco de dados.
    \item Gerar rotas HTTP automaticamente.
\end{enumerate}

\vspace{0.2cm}
\textbf{3.} Sempre que definimos um método interno dentro de uma Classe em Python, precisamos passar a palavra \texttt{self} como o primeiro parâmetro obrigatório. O que o \texttt{self} representa no código?
\begin{enumerate}[label=(\alph*)]
    \item A variável global que armazena senhas.
    \item \correto{A própria instância do objeto, permitindo que a classe enxergue e altere seus próprios atributos internos.}
    \item O ponteiro de memória do servidor.
    \item Uma diretiva de segurança de acesso.
\end{enumerate}

\vspace{0.2cm}
\textbf{4.} O que faz "pip install" pacote?
\begin{enumerate}[label=(\alph*)]
    \item Instala um objeto no código.
    \item \correto{Baixa alguma biblioteca para o python.}
    \item Gera um component no Ionic.
    \item Cria uma rota http para acionamento do Flask.
\end{enumerate}

\vspace{0.2cm}
\textbf{5.} Onde é utilizado linguagens de programação como o Python para gerenciar o sistema?
\begin{enumerate}[label=(\alph*)]
    \item Frontend
    \item \correto{API}
    \item Banco de Dados
    \item Sistema operacional
\end{enumerate}

\vspace{0.2cm}
\textbf{6.} O Docker revolucionou o modo como distribuímos software. Sobre os conceitos fundamentais desta tecnologia, assinale a diferença correta entre "Imagem" e "Container":
\begin{enumerate}[label=(\alph*)]
    \item Containers rodam apenas no sistema Windows, enquanto as Imagens são feitas para Linux.
    \item Não há diferença prática, ambos são nomes diferentes que a comunidade usa.
    \item \correto{A Imagem é uma "planta-baixa" estática (pacote) com os arquivos. O Container é a Imagem instanciada, executando como processo.}
    \item A Imagem é o banco de dados, e o container é apenas o Frontend.
\end{enumerate}

\vspace{0.2cm}
\textbf{7. [V/F]} ( ) O comando \texttt{git pull} é utilizado no terminal local para baixar as atualizações do servidor remoto (GitHub) e sincronizá-las com a pasta física.
\justificativa{Verdadeiro.}

\vspace{0.2cm}
\textbf{8.} De forma objetiva, o que é uma API?
\begin{enumerate}[label=(\alph*)]
    \item Apenas a interface gráfica de um site.
    \item \correto{Um programa conectado na internet que funciona quando chamado por outro sistema.}
    \item O mesmo que um sistema operacional completo.
    \item Um tipo exclusivo de banco de dados relacional.
\end{enumerate}

\vspace{0.2cm}
\textbf{9. [V/F]} ( ) Se uma API processa uma requisição perfeitamente, encontra o usuário no banco de dados e devolve os dados para a tela com sucesso, o \textit{Status Code} de resposta será o \textbf{404 Not Found}.
\justificativa{Falso. O código 404 indica erro de recurso não encontrado. Para sucesso, devolve-se 200 OK.}

\vspace{0.2cm}
\textbf{10. [V/F]} ( ) O método \texttt{GET} é considerado seguro para enviar dados sensíveis, como senhas de login, porque seus parâmetros são encapsulados no corpo (\textit{Body}) da requisição, invisíveis na URL.
\justificativa{Falso. O método GET expõe os parâmetros na URL. O método POST é o mais adequado.}

\vspace{0.2cm}
\textbf{11. [V/F]} ( ) No comando de subir um servidor web no Docker, a \textit{flag} \texttt{-p 8080:80} significa que estamos criando um túnel que pega acessos externos da porta 8080 da sua máquina real e injeta na porta 80 dentro do container.
\justificativa{Verdadeiro.}

\vspace{0.2cm}
\textbf{12.} O que representa um .git num projeto?
\begin{enumerate}[label=(\alph*)]
    \item Uma cópia completa do repositório em outro servidor
    \item Um snapshot registrado do estado dos arquivos
    \item Apenas a diferença entre duas branches
    \item \correto{Um arquivo/pasta de configuração do repositório}
\end{enumerate}

\vspace{0.2cm}
\textbf{13. [Sequência]} Relacione a \textbf{Coluna A} (Método HTTP) com a \textbf{B} (Ação CRUD).
(1) \texttt{GET}, (2) \texttt{POST}, (3) \texttt{PUT}, (4) \texttt{DELETE}.\\
\textbf{Coluna B:}
\begin{itemize}[leftmargin=*, label={}]
    \item ( \resp{2} ) \textbf{C}reate: Inserir um novo recurso.
    \item ( \resp{1} ) \textbf{R}ead: Ler e visualizar informações.
    \item ( \resp{3} ) \textbf{U}pdate: Substituir ou editar dados.
    \item ( \resp{4} ) \textbf{D}elete: Excluir um registro.
\end{itemize}

\vspace{0.2cm}
\textbf{14. [Sequência]} Relacione \textbf{A} (Comando) com \textbf{B} (Ação).
(1) \texttt{run}, (2) \texttt{ps}, (3) \texttt{exec}, (4) \texttt{stop}.\\
\textbf{Coluna B:}
\begin{itemize}[leftmargin=*, label={}]
    \item ( \resp{2} ) Lista containers em execução.
    \item ( \resp{3} ) Entra num container funcionando.
    \item ( \resp{1} ) Baixa imagem e dá a partida.
    \item ( \resp{4} ) Interrompe e desliga o container.
\end{itemize}

\end{multicols}

\vspace{0.3cm}
\hrule
\vspace{0.3cm}
\textbf{15.} Analise o trecho de código abaixo de uma API Flask:
\begin{verbatim}
tarefas = [ {"id": 1, "titulo": "Estudar Docker"}, {"id": 2, "titulo": "Aprender Flask"} ]

@app.route('/tarefas', methods=['GET'])
def listar_tarefas():
    return jsonify(tarefas)
\end{verbatim}
Se um aluno enviar uma requisição \texttt{GET} para \texttt{http://localhost:5000/tarefas}, o que será devolvido na tela?
\begin{enumerate}[label=(\alph*)]
    \item Uma página HTML contendo o título das tarefas.
    \item O código dará erro 404 pois a porta está incorreta.
    \item \correto{Um arquivo JSON com a lista de tarefas e o status 200 OK.}
    \item Um texto puro dizendo "Estudar Docker" e "Aprender Flask".
\end{enumerate}

\vfill
\begin{center}
    \textit{Sucesso na atividade! Atenciosamente Prof. André Moraes}
\end{center}

\end{document}
